\chapter{Typed Morphisms, Spans, and Polyxan Rewriting}
\label{ch:polyxan-rewriting}

Chapter~\ref{ch:polyxan-basic} introduced Polyxan hypergraphs as typed,
tagged, sheaf-valued, curvature-bearing combinatorial structures.  This chapter
develops their morphisms, spans, rewriting systems, and reduction theory.
These provide the algebraic machinery underlying the Polycompiler and the
media-quine attractor of Part~II.

% ===============================================================
\section{Typed Morphisms Revisited}
% ===============================================================

We recall that a Polyxan morphism $f : G \to H$ is type-preserving, curvature-preserving, and induces a morphism of sheaves on the incidence complexes.  We now refine this definition to include interface embeddings essential for rewriting.

\begin{definition}[Typed mono]
A morphism $m : K \hookrightarrow G$ in $\Polyx$ is a \emph{typed mono} if
$m$ is injective on vertices and hyperedges and reflects types:
\[
t(m(a)) = t(a), \qquad \lambda(m(e)) = \lambda(e).
\]
Typed monos serve as \emph{interfaces} for rewrite rules.
\end{definition}

Typed monos correspond to semantic substructures that can be matched,
extracted, or replaced.

% ===============================================================
\section{Spans, Cospans, and the Double Category PolyxSpan}
% ===============================================================

\subsection{Spans of Polyxan hypergraphs}

\begin{definition}[Span]
A \emph{span} in $\Polyx$ is a diagram
\[
L \xleftarrow{\;\ell\;} K \xrightarrow{\;r\;} R,
\]
where $\ell,r$ are typed monos.  $K$ is the \emph{interface}.
\end{definition}

$K$ represents the part of $L$ to be rewritten, and $R$ the replacement.

\subsection{Cospans}

Dually:

\begin{definition}[Cospan]
A \emph{cospan} is a diagram
\[
L \xrightarrow{\;\ell'\;} C \xleftarrow{\;r'\;} R,
\]
with typed monos $\ell',r'$.
\end{definition}

Cospans will appear later when we interpret Polyxan rewrites as semantic
\emph{coalescence} or structural \emph{fusion} operations.

\subsection{Double category}

\begin{theorem}
The data of Polyxan hypergraphs, typed morphisms, spans, and cospans forms a
strict double category
\[
\PolyxSpan,
\]
whose objects are Polyxan hypergraphs, horizontal 1-morphisms are spans,
vertical 1-morphisms are typed morphisms, and 2-cells are commuting diagrams
of interface-preserving maps.
\end{theorem}

\begin{proof}[Sketch]
Horizontal composition of spans uses fibered products
\[
(L \xleftarrow{\ell} K \xrightarrow{r} R)
\circ
(R \xleftarrow{\ell'} K' \xrightarrow{r'} S)
:=
L \xleftarrow{\ell\circ\pi_1} K \times_R K' \xrightarrow{r'\circ \pi_2} S
\]
and vertical composition uses ordinary composition of typed morphisms.
Compatibility of sources and targets follows from preservation of type and
curvature.  The interchange law holds by construction of fibered products in
the presheaf topos equivalent to $\Polyx$ (Theorem~\ref{thm:polyx-topos}).
\end{proof}

PolyxSpan is the categorical home of rewriting.

% ===============================================================
\section{Typed Rewrite Rules}
% ===============================================================

\begin{definition}[Rewrite rule]
A \emph{typed rewrite rule} is a span
\[
L \xleftarrow{\;\ell\;} K \xrightarrow{\;r\;} R,
\]
where:
\begin{enumerate}
  \item $L$ is the \emph{left-hand side};
  \item $R$ is the \emph{right-hand side};
  \item $K$ is the \emph{interface}, embedded into $L$ and $R$ via typed monos.
\end{enumerate}
\end{definition}

Given an embedding $m : L \to G$, a match of $L$ in $G$, a rewrite replaces
$m(L)$ by $m(K)$ glued to $R$, provided sheaf and curvature constraints hold.

\begin{definition}[Polyxan rewriting step]
A rewriting step is a pushout in $\Polyx$:
\[
\begin{tikzcd}
K \ar[r,"r"] \ar[d,"\ell"'] & R \ar[d] \\
L \ar[r] & G'
\end{tikzcd}
\]
where $G'$ is the hypergraph obtained by replacing $L$ with $R$ along $K$.
\end{definition}

% ===============================================================
\section{Sesquicategory of Polyxan Rewriting Systems}
% ===============================================================

Rewrite rules assemble into a sesquicategory (2-category without identity
2-cells).

\begin{definition}[Sesquicategory Rewrite(Polyx)]
Objects: Polyxan hypergraphs.  
1-morphisms: finite sequences of rewrite steps.  
2-cells: commutation diagrams between rewrite sequences induced by interface
maps.
\end{definition}

\begin{theorem}[Church–Rosser Lemma]
If a rewrite system is locally confluent and terminating, then it is confluent:
whenever $G \to^\ast H_1$ and $G \to^\ast H_2$, there exists $J$ with
$H_1 \to^\ast J$ and $H_2 \to^\ast J$.
\end{theorem}

\begin{theorem}[Newman's Lemma]
If a rewrite system is terminating and locally confluent, it is globally
confluent.
\end{theorem}

Both will hold for the reduction systems defined below.

% ===============================================================
\section{The Extraction–Hoisting–Reduction Cycle}
% ===============================================================

Rewriting in Polyx is driven by a canonical cycle of operations.

\subsection{Extraction}

Given $G$, extraction removes syntactically redundant or curvature-neutral
substructures, producing a simplified $G_{\mathrm{ext}}$.  
Formally, extraction is a functor
\[
\mathsf{Ext}: \Polyx \to \Polyx,
\]
defined by eliminating subobjects whose sheaf sections are contractible and
whose curvature contribution is zero.

\subsection{Hoisting}

Hoisting reorganizes hyperedges by raising low-level representations to more
abstract equivalents.  It is a functor
\[
\mathsf{Hoist}: \Polyx \to \Polyx,
\]
preserving sheaf data and hyperedge types, but collapsing unnecessary modality
distinctions.

\subsection{Reduction}

Reduction applies curvature-decreasing rewrite rules.  
Formally,
\[
\mathsf{Red}: \Polyx \to \Polyx
\]
is defined by repeatedly applying a finite set of rewrite rules
$\{L_i \leftarrow K_i \rightarrow R_i\}$ chosen to strictly reduce a curvature
measure.

\subsection{Canonical endofunctor}

\begin{definition}[EHR endofunctor]
The \emph{extraction–hoisting–reduction cycle} is the endofunctor
\[
\mathsf{EHR} := \mathsf{Red} \circ \mathsf{Hoist} \circ \mathsf{Ext}
: \Polyx \to \Polyx.
\]
\end{definition}

EHR governs Polyxan evolution and Polycompiler convergence.

% ===============================================================
\section{Termination via Curvature Decrease}
% ===============================================================

The curvature assignment $\kappa_G$ defines a height function on $\Polyx$.

\begin{theorem}[Curvature-decreasing termination]
If every reduction rule strictly decreases at least one curvature invariant
(hyperedge deficit, sectional curvature, sheaf curvature, or nerve Euler
characteristic), then all reduction sequences terminate.
\end{theorem}

\begin{proof}[Sketch]
The curvature tuple
\[
\vec{\kappa}(G) := (\sum_e \delta(e),\; \sum_x |K_x|,\; \sum_{U,V,W} \Omega(U,V,W),\; \chi(G))
\]
takes values in $\mathbb{N}^4$ with lexicographic order.  
Reduction strictly decreases $\vec{\kappa}(G)$, hence no infinite descending
chain exists.
\end{proof}

This is the discrete analogue of the Lyapunov monotonicity established for
RSVP fields.

% ===============================================================
\section{Polyxan Energy and Media-Quine Critical Points}
% ===============================================================

\begin{definition}[Polyxan energy]
Define
\[
\mathcal{E}[G] :=
\sum_{e\in E(G)} \delta(e)
\;+\;
\sum_{x\in G} |K_x|
\;+\;
\sum_{U,V,W} \Omega(U,V,W)
\;+\;
|\chi(G)|,
\]
a scalar energy functional on $\Polyx$.
\end{definition}

\begin{theorem}
$\mathsf{EHR}$ decreases $\mathcal{E}[G]$ strictly unless $G$ is a media quine.
\end{theorem}

Thus media quines are precisely the critical points of $\mathcal{E}$.

% ===============================================================
\section{Unique Normal Forms}
% ===============================================================

\begin{theorem}[Polyxan Unique Normal Form Theorem]
For every finite Polyxan hypergraph $G$, the reduction system defined by
$\mathsf{EHR}$ is terminating and confluent.  
Thus $G$ has a unique normal form $G^\ast$ satisfying
\[
G \;\to^\ast\; G^\ast,
\qquad
\mathsf{EHR}(G^\ast) = G^\ast.
\]
Moreover, $G^\ast$ is a media quine.
\end{theorem}

\begin{proof}[Sketch]
Termination follows from curvature decrease.  
Local confluence is verified by checking overlaps of extraction, hoisting, and
reduction rules.  
Newman's Lemma then implies global confluence.  
Thus every $G$ reduces uniquely to a fixed point $G^\ast$, which must satisfy
$\mathsf{EHR}(G^\ast)=G^\ast$, hence is a media quine.
\end{proof}

% ===============================================================
\section{Summary}
% ===============================================================}

This chapter developed the algebraic and rewrite-theoretic machinery of Polyxan
hypergraphs:
\begin{itemize}
  \item spans and cospans forming the double category $\PolyxSpan$;
  \item typed rewrite rules via $L\leftarrow K\to R$;
  \item the sesquicategory of rewriting systems and confluence theorems;
  \item the extraction--hoisting--reduction cycle as an endofunctor on $\Polyx$;
  \item curvature-decreasing termination and the Polyxan energy functional;
  \item the Polyxan Unique Normal Form Theorem identifying media quines as the
        canonical fixed points.
\end{itemize}

These structures constitute the algebraic engine that powers the entire
Polycompiler architecture, and they prepare the ground for the spectral,
sheaf-theoretic, and categorical analyses of the chapters that follow.
